\documentclass[11pt]{article}
\usepackage[margin=1in]{geometry}
\usepackage{amsmath,amssymb,amsthm,mathtools,bm}
\usepackage{booktabs}

\title{Formalization of Constrained Flow Matching v1}
\author{Project Notes}
\date{April 27, 2026}

\begin{document}
\maketitle

\section{Goal and Setting}
We want to learn a velocity field that transports $p_0$ to $p_1$, while enforcing intermediate-time moment constraints.
In v1, the state is $x \in \mathbb{R}^2$.

\paragraph{Given}
\begin{itemize}
    \item Initial and terminal marginals $p_0, p_1$.
    \item A coupling over endpoint pairs $(x_0, x_1)$ (configured as exact discrete OT or random permutation per batch).
    \item Constraint times $\mathcal{T}_c = \{0.25, 0.5, 0.75\}$.
    \item Target moments $m_k$ at each $t_k \in \mathcal{T}_c$.
\end{itemize}

\paragraph{Wanted}
Learn:
\begin{itemize}
    \item a path-correction model $g_\theta(t,x_0,x_1)$,
    \item a velocity model $v_\psi(t,x)$,
\end{itemize}
such that induced paths satisfy moment constraints and $v_\psi$ matches path velocities.

\section{Reference: Standard CFM}
Standard conditional flow matching uses a fixed interpolation path $x_t$ (often linear) and target velocity $u_t = \partial_t x_t$, then trains
\[
\mathcal{L}_{\mathrm{CFM}}(\psi)
=
\mathbb{E}_{(x_0,x_1),\,t}
\left[
\left\|
v_\psi(t,x_t) - u_t
\right\|^2
\right].
\]

\section{Path Family Used in This Project}
\subsection{Baseline Path}
\[
\ell_t(x_0,x_1) = (1-t)x_0 + tx_1.
\]

\subsection{Constrained-Correction Path}
We use an endpoint-preserving correction:
\[
x_t = \ell_t + a(t)\, g_\theta(t,x_0,x_1),
\qquad
a(t)=t(1-t).
\]
Because $a(0)=a(1)=0$, endpoints are exactly preserved: $x_{t=0}=x_0$, $x_{t=1}=x_1$.

\subsection{Target Velocity for CFM in Constrained Mode}
\[
u_t = \partial_t x_t
=
(x_1-x_0)
+ (1-2t)\,g_\theta(t,x_0,x_1)
+ a(t)\,\partial_t g_\theta(t,x_0,x_1).
\]
This is exactly what is computed in code using autograd for $\partial_t g_\theta$.

\section{Moment Constraints}
\subsection{Feature Map Used in v1}
For a batch $\{x^{(i)}\}_{i=1}^B \subset \mathbb{R}^2$:
\[
\hat{\mu} = \frac{1}{B}\sum_{i=1}^B x^{(i)},
\qquad
\hat{\Sigma} = \frac{1}{B}\sum_{i=1}^B
\big(x^{(i)}-\hat{\mu}\big)\big(x^{(i)}-\hat{\mu}\big)^\top.
\]
Define
\[
\phi(\{x^{(i)}\}_{i=1}^B)
=
\begin{bmatrix}
\hat{\mu} \\
\mathrm{vec}(\hat{\Sigma})
\end{bmatrix}
\in \mathbb{R}^{6}.
\]

\subsection{Residual at Constraint Times}
For each $t_k \in \mathcal{T}_c$:
\[
c_k(\theta) = \phi\!\left(\{x_{t_k}^{(i)}(\theta)\}_{i=1}^B\right) - m_k.
\]
Small $\|c_k\|$ means constraints are well satisfied.

\section{How Target Moments $m_k$ are Built in v1}
In the benchmark, $(p_0,p_1)$ are Gaussians:
\[
p_0=\mathcal{N}(\mu_0,\Sigma_0), \qquad p_1=\mathcal{N}(\mu_1,\Sigma_1).
\]
Define the Gaussian OT map matrix:
\[
A
=
\Sigma_0^{-1/2}
\left(
\Sigma_0^{1/2}\Sigma_1\Sigma_0^{1/2}
\right)^{1/2}
\Sigma_0^{-1/2}.
\]
Define
\[
B_t=(1-t)I+tA,
\qquad
s(t)=1+\kappa\,t(1-t).
\]
Then target mean/covariance are:
\[
\mu_t = (1-t)\mu_0 + t\mu_1,
\qquad
\Sigma_t = s(t)^2\, B_t \Sigma_0 B_t^\top.
\]
Finally:
\[
m_k
=
\begin{bmatrix}
\mu_{t_k}\\
\mathrm{vec}(\Sigma_{t_k})
\end{bmatrix}.
\]

\section{Loss Terms Used in the Current Implementation}
\subsection{Path Regularizer}
\[
R(\theta)
=
\alpha\,
\mathbb{E}\!\left[\|u_t-(x_1-x_0)\|^2\right]
+\beta\,
\mathbb{E}\!\left[\|\partial_t u_t\|^2\right].
\]

\paragraph{Interpretation}
\begin{itemize}
    \item First term keeps the learned path velocity near linear transport velocity.
    \item Second term penalizes rapid temporal variation of velocity (temporal smoothness).
\end{itemize}

\subsection{Augmented Lagrangian for Path Learning (Stage A/C)}
\[
\mathcal{L}_g(\theta,\lambda)
=
R(\theta)
+ \sum_{k}\lambda_k^\top c_k(\theta)
+ \frac{\rho}{2}\sum_k \|c_k(\theta)\|^2.
\]
Multiplier updates (with clipping in code):
\[
\lambda_k \leftarrow
\mathrm{clip}\!\left(\lambda_k + \rho\,c_k(\theta),\,-\lambda_{\max},\,\lambda_{\max}\right).
\]

\subsubsection*{Algorithm 1: Mini-batch augmented-Lagrangian updates (code-equivalent)}
\textbf{Inputs:} constraint times $\mathcal{T}_c$, target features $\{m_k\}_{t_k\in\mathcal{T}_c}$, hyperparameters $(\alpha,\beta,\rho,\eta_{\mathrm{joint}},\lambda_{\max})$, and step counts $(N_A,N_B,N_C)$.

\begin{enumerate}
    \item Initialize multipliers: $\lambda_k \gets 0$ for all $t_k\in\mathcal{T}_c$.
    \item \textbf{Stage A (optimize $\theta$ only):} for $s=1,\dots,N_A$:
    \begin{enumerate}
        \item Sample paired mini-batch $\{(x_0^{(i)},x_1^{(i)})\}_{i=1}^B$.
        \item Sample $t^{(i)}\sim\mathrm{Uniform}(0,1)$ and compute constrained velocity targets $u_{t^{(i)}}(\theta)$.
        \item Compute regularizer
        \[
        R(\theta)\gets
        \alpha\,\frac{1}{B}\sum_i\|u_{t^{(i)}}-(x_1^{(i)}-x_0^{(i)})\|_2^2
        +\beta\,\frac{1}{B}\sum_i\|\partial_t u_{t^{(i)}}\|_2^2.
        \]
        \item For each $t_k\in\mathcal{T}_c$, compute
        \[
        x_{t_k}^{(i)}\gets
        (1-t_k)x_0^{(i)}+t_kx_1^{(i)}+t_k(1-t_k)\,g_\theta(t_k,x_0^{(i)},x_1^{(i)}),
        \]
        \[
        \hat{\mu}_k\gets \frac{1}{B}\sum_i x_{t_k}^{(i)},\qquad
        \hat{\Sigma}_k\gets \frac{1}{B}\sum_i(x_{t_k}^{(i)}-\hat{\mu}_k)(x_{t_k}^{(i)}-\hat{\mu}_k)^\top,
        \]
        \[
        c_k(\theta)\gets
        \begin{bmatrix}
        \hat{\mu}_k\\
        \mathrm{vec}(\hat{\Sigma}_k)
        \end{bmatrix}
        -m_k.
        \]
        \item Assemble AL objective
        \[
        \mathcal{L}_g\gets
        R(\theta)+\sum_k\lambda_k^\top c_k(\theta)+\frac{\rho}{2}\sum_k\|c_k(\theta)\|_2^2.
        \]
        \item Apply Adam step on $\theta$ with gradient $\nabla_\theta\mathcal{L}_g$.
        \item Update multipliers for all $t_k\in\mathcal{T}_c$:
        \[
        \lambda_k\gets
        \mathrm{clip}\!\left(\lambda_k+\rho\,c_k(\theta),-\lambda_{\max},\lambda_{\max}\right).
        \]
    \end{enumerate}
    \item \textbf{Stage B (freeze $\theta$, optimize $\psi$):} for $s=1,\dots,N_B$, sample a mini-batch and apply Adam step on $\psi$ with $\nabla_\psi\mathcal{L}_{\mathrm{CFM}}$.
    \item \textbf{Stage C (joint optimize $(\psi,\theta)$):} for $s=1,\dots,N_C$:
    \begin{enumerate}
        \item Sample a mini-batch and recompute $\mathcal{L}_g(\theta,\lambda)$ exactly as in Stage A.
        \item Compute $\mathcal{L}_{\mathrm{CFM}}(\psi;\theta)$ and
        \[
        \mathcal{L}_{\mathrm{joint}}\gets \mathcal{L}_{\mathrm{CFM}}+\eta_{\mathrm{joint}}\mathcal{L}_g.
        \]
        \item Apply Adam step on $(\psi,\theta)$ with gradient $\nabla_{(\psi,\theta)}\mathcal{L}_{\mathrm{joint}}$.
        \item Update all multipliers with the same clipped rule:
        \[
        \lambda_k\gets
        \mathrm{clip}\!\left(\lambda_k+\rho\,c_k(\theta),-\lambda_{\max},\lambda_{\max}\right).
        \]
    \end{enumerate}
\end{enumerate}

\subsection{CFM Loss for Velocity Network}
\[
\mathcal{L}_{\mathrm{CFM}}(\psi;\theta)
=
\mathbb{E}_{(x_0,x_1),t}
\left[
\left\|
v_\psi(t,x_t(\theta)) - u_t(\theta)
\right\|^2
\right].
\]

\subsection{Joint Finetune Objective (Stage C)}
\[
\mathcal{L}_{\mathrm{joint}}(\psi,\theta,\lambda)
=
\mathcal{L}_{\mathrm{CFM}}(\psi;\theta)
+ \eta_{\mathrm{joint}}\,\mathcal{L}_g(\theta,\lambda).
\]

\section{Current Optimization Loop (Exactly What Code Does)}
\subsection*{Sampling per Step}
At each step:
\begin{enumerate}
    \item Sample endpoint mini-batch from $p_0$ and $p_1$.
    \item Solve exact discrete OT assignment on this batch to obtain paired $(x_0^{(i)},x_1^{(i)})$.
    \item Sample $t \sim \mathrm{Uniform}(0,1)$ for CFM/regularizer terms.
\end{enumerate}

\subsection*{Stage A: Pretrain Path Model $g_\theta$}
For $N_A$ steps:
\begin{enumerate}
    \item Compute $\mathcal{L}_g(\theta,\lambda)$.
    \item Gradient step on $\theta$ (Adam, lr\_g).
    \item Update $\lambda_k \leftarrow \lambda_k + \rho c_k$ (with clipping).
\end{enumerate}

\subsection*{Stage B: Train Velocity Model $v_\psi$ with Frozen $g_\theta$}
For $N_B$ steps:
\begin{enumerate}
    \item Freeze $\theta$.
    \item Compute $\mathcal{L}_{\mathrm{CFM}}(\psi;\theta)$.
    \item Gradient step on $\psi$ (Adam, lr\_v).
\end{enumerate}

\subsection*{Stage C: Joint Finetune}
For $N_C$ steps:
\begin{enumerate}
    \item Unfreeze both $\theta,\psi$.
    \item Compute $\mathcal{L}_{\mathrm{joint}}$.
    \item Gradient step on $(\theta,\psi)$.
    \item Update multipliers $\lambda_k$ again.
\end{enumerate}

\subsection*{Baseline Mode}
In baseline mode there is no $g_\theta$:
\[
x_t=\ell_t,\qquad u_t=x_1-x_0.
\]
Only Stage B is optimized.

\section{Why This is More Complex than ``Path First, then Flow Matching''}
Your suggested two-step idea is:
\begin{enumerate}
    \item Learn paths (e.g., with a smoothness/Lipschitz regularizer),
    \item then train flow matching on those paths.
\end{enumerate}

Current implementation adds two extra mechanisms:
\begin{itemize}
    \item \textbf{Augmented Lagrangian} instead of fixed penalties, to push constraints more aggressively.
    \item \textbf{Stage C joint finetune}, to co-adapt path model and velocity model after separate training.
\end{itemize}

So, yes: the present loop is strictly more complicated than a plain two-stage approach.
It can improve constraint handling, but it also introduces additional instability/tuning sensitivity.

\section{How Reported Metrics Are Computed}
This section defines the fields written in \texttt{metrics.json} and \texttt{comparison.json}.

\subsection*{Constraint metrics}
For each constrained time $t_k \in \mathcal{T}_c$, compute residual vector
\[
c_k = \phi\!\left(\{x_{t_k}^{(i)}\}_{i=1}^B\right)-m_k.
\]
Then:
\[
\texttt{constraint\_residual\_norms}[t_k] = \|c_k\|_2.
\]
The summary scalar is:
\[
\texttt{constraint\_residual\_avg}
=
\frac{1}{|\mathcal{T}_c|}
\sum_{t_k\in\mathcal{T}_c}\|c_k\|_2.
\]

\subsection*{Validation CFM metric}
On an evaluation batch:
\begin{enumerate}
    \item sample paired $(x_0,x_1)$ via the configured coupling mode,
    \item sample $t \sim \mathrm{Uniform}(0,1)$,
    \item build $(x_t,u_t)$ from current mode (baseline or constrained),
    \item evaluate
    \[
    \texttt{cfm\_val\_loss}
    =
    \frac{1}{B}\sum_{i=1}^B
    \left\|
    v_\psi(t_i,x_{t_i}^{(i)})-u_{t_i}^{(i)}
    \right\|_2^2.
    \]
\end{enumerate}

\subsection*{Intermediate-time Wasserstein metric}
For each constrained time $t_k \in \mathcal{T}_c$:
\begin{enumerate}
    \item Sample $x_0^{(i)} \sim p_0$, $i=1,\dots,N$.
    \item Integrate the learned velocity field from $0$ to $t_k$ with Euler to obtain $\hat{x}_{t_k}^{(i)}$.
    \item Compute empirical generated moments:
    \[
    \hat{\mu}_{t_k} = \frac{1}{N}\sum_{i=1}^N \hat{x}_{t_k}^{(i)},
    \qquad
    \hat{\Sigma}_{t_k} = \frac{1}{N}\sum_{i=1}^N
    (\hat{x}_{t_k}^{(i)}-\hat{\mu}_{t_k})(\hat{x}_{t_k}^{(i)}-\hat{\mu}_{t_k})^\top.
    \]
    \item Let target moments be $(\mu_{t_k},\Sigma_{t_k})$ from the analytic Gaussian bridge.
\end{enumerate}
The reported per-time metric uses Gaussian $W_2$ distance:
\[
\texttt{intermediate\_w2\_gaussian}[t_k]
=
W_2\!\left(
\mathcal{N}(\hat{\mu}_{t_k},\hat{\Sigma}_{t_k}),
\mathcal{N}(\mu_{t_k},\Sigma_{t_k})
\right),
\]
with closed-form
\[
W_2^2
=
\|\hat{\mu}_{t_k}-\mu_{t_k}\|_2^2
+
\mathrm{Tr}\!\left(
\hat{\Sigma}_{t_k}
+
\Sigma_{t_k}
-
2\left(\Sigma_{t_k}^{1/2}\hat{\Sigma}_{t_k}\Sigma_{t_k}^{1/2}\right)^{1/2}
\right).
\]
The summary scalar is:
\[
\texttt{intermediate\_w2\_gaussian\_avg}
=
\frac{1}{|\mathcal{T}_c|}\sum_{t_k\in\mathcal{T}_c}
\texttt{intermediate\_w2\_gaussian}[t_k].
\]

\paragraph{Sample-based empirical OT variant}
To support non-Gaussian settings, we also report a sample-based empirical $W_2$ proxy:
\begin{enumerate}
    \item Generate trajectory samples $\{\hat{x}_{t_k}^{(i)}\}_{i=1}^N$ by integrating $v_\psi$ from $0$ to $t_k$.
    \item Draw target samples $\{y_{t_k}^{(j)}\}_{j=1}^N$ from the reference distribution at $t_k$.
    \item Compute the optimal bipartite matching $\pi^\star$ minimizing squared Euclidean cost:
    \[
    \pi^\star = \arg\min_{\pi \in S_N}\sum_{i=1}^N
    \left\|\hat{x}_{t_k}^{(i)}-y_{t_k}^{(\pi(i))}\right\|_2^2.
    \]
\end{enumerate}
The reported per-time empirical metric is:
\[
\texttt{intermediate\_empirical\_w2}[t_k]
=
\sqrt{
\frac{1}{N}
\sum_{i=1}^N
\left\|\hat{x}_{t_k}^{(i)}-y_{t_k}^{(\pi^\star(i))}\right\|_2^2
}.
\]
and the average is:
\[
\texttt{intermediate\_empirical\_w2\_avg}
=
\frac{1}{|\mathcal{T}_c|}\sum_{t_k\in\mathcal{T}_c}
\texttt{intermediate\_empirical\_w2}[t_k].
\]
In code this uses exact discrete OT matching (Hungarian algorithm), so evaluation sample count is configured separately to control CPU cost.

\subsection*{Path energy proxy}
Using the same evaluation batch/time samples as \texttt{cfm\_val\_loss}:
\[
\texttt{path\_energy\_proxy}
=
\frac{1}{B}\sum_{i=1}^B \|u_{t_i}^{(i)}\|_2^2.
\]
This is a diagnostic quantity, not an objective minimized directly at evaluation time.

\subsection*{Transport quality metrics}
To measure endpoint transport induced by $v_\psi$:
\begin{enumerate}
    \item Sample $x_0^{(i)} \sim p_0$, $i=1,\dots,N$.
    \item Integrate forward with explicit Euler over $n_{\text{steps}}$ steps:
    \[
    x^{(i)}_{n+1}
    =
    x^{(i)}_{n}
    +
    \Delta t\, v_\psi(t_n, x^{(i)}_n),
    \quad
    \Delta t = \frac{1}{n_{\text{steps}}},
    \quad
    t_n = n\Delta t.
    \]
    \item Let $\hat{x}_1^{(i)} = x^{(i)}_{n_{\text{steps}}}$.
    \item Compute empirical endpoint statistics:
    \[
    \hat{\mu}_1 = \frac{1}{N}\sum_{i=1}^N \hat{x}_1^{(i)},
    \qquad
    \hat{\Sigma}_1=\frac{1}{N}\sum_{i=1}^N
    (\hat{x}_1^{(i)}-\hat{\mu}_1)(\hat{x}_1^{(i)}-\hat{\mu}_1)^\top.
    \]
\end{enumerate}
Then:
\[
\texttt{transport\_mean\_error\_l2}
=
\|\hat{\mu}_1-\mu_1\|_2,
\]
\[
\texttt{transport\_cov\_error\_fro}
=
\|\hat{\Sigma}_1-\Sigma_1\|_F,
\]
\[
\texttt{transport\_score}
=
\texttt{transport\_mean\_error\_l2}
+
\texttt{transport\_cov\_error\_fro}.
\]

\subsection*{Stage-history diagnostics}
During training, the history log also records:
\begin{itemize}
    \item \texttt{avg\_residual\_norm}: mini-batch average of $\|c_k\|_2$ over constrained times,
    \item \texttt{loss}: stage objective value,
    \item and stage-specific terms (e.g., regularizer or \texttt{cfm\_loss} in Stage C).
\end{itemize}

\section{Example Experiment Settings (Current Repository)}
This section summarizes the experiment families currently implemented and logged in this repository.

\subsection*{A. Primary trained benchmark: 2D Gaussian $\rightarrow$ 2D Gaussian}
\begin{itemize}
    \item \textbf{Data family}: both endpoints are Gaussians in $\mathbb{R}^2$.
    \item \textbf{Targets}: intermediate moment targets are analytic Gaussian bridge moments at
    $t\in\{0.25,0.50,0.75\}$.
    \item \textbf{Pair coupling options}: exact discrete OT (\texttt{data=gaussian\_ot}) or random permutation pairing (\texttt{data=gaussian\_random}).
    \item \textbf{Training modes}: baseline (linear path), constrained (Stage A/B/C), and comparison (run both).
    \item \textbf{Typical profiles}: \texttt{train=default}, \texttt{train=ab\_only} (Stage C disabled), and \texttt{train=smoke}.
    \item \textbf{Representative output roots}: \texttt{outputs/hparam\_sweep/}, \texttt{outputs/ab\_only\_sweep/}, and timestamped comparison folders under \texttt{outputs/<date>/...}.
\end{itemize}

\subsection*{B. Bridge-SDE preview: Gaussian start with bridge-like stochastic dynamics}
\begin{itemize}
    \item \textbf{Purpose}: geometry/diagnostic preview prior to training on more complex non-Gaussian trajectories.
    \item \textbf{Initial law}: Gaussian $p_0=\mathcal{N}(\mu_0,\Sigma_0)$ from bridge-preview parameters.
    \item \textbf{Simulated dynamics} (from \texttt{bridge\_sde.py}):
    \[
    dX_t = v_x\,dt + \sigma_x\,dW_t^{(x)},
    \]
    \[
    dY_t =
    -\kappa_b\,g(X_t)\,Y_t\,dt
    + \sigma_y\!\left(1-\delta_b\,g(X_t)\right)dW_t^{(y)},
    \qquad
    g(x)=\exp\!\left(-\frac{(x-c)^2}{2w^2}\right).
    \]
    \item \textbf{Current status}: used in notebook/preview artifacts (snapshots, spread plots, animation), not yet integrated as a full training target in the main Stage A/B/C experiment pipeline.
    \item \textbf{Output root}: \texttt{outputs/preview\_bridge\_sde/<timestamp>/}.
\end{itemize}

\paragraph{Scope note}
At present, the optimization formalized in Sections 1--9 corresponds to the Gaussian-to-Gaussian benchmark pipeline.
The Bridge-SDE path is currently a validated simulation/visualization track used to prepare future non-Gaussian training settings.

\section{Practical Hyperparameters in the Loop}
\begin{center}
\begin{tabular}{ll}
\toprule
Symbol / Config & Role \\
\midrule
$\alpha$ & weight of velocity-deviation regularizer \\
$\beta$ & weight of temporal smoothness term \\
$\rho$ & augmented-Lagrangian quadratic and multiplier update scale \\
$\eta_{\mathrm{joint}}$ & weight of $\mathcal{L}_g$ inside Stage C objective \\
lr\_g, lr\_v & learning rates for path and velocity networks \\
$N_A,N_B,N_C$ & number of steps in each stage \\
\bottomrule
\end{tabular}
\end{center}

\end{document}
